\renewcommand{\Anno}{1997}

\renewcommand{\NumeroProblema}{4}

\renewcommand{\Testo}{
Un \emph{trisecatore} è uno strumento che, dato un qualunque segmento, traccia due punti che lo suddividono in 3 parti uguali.

Dimostrare che, con il solo aiuto di una riga ed un trisecatore, si riesce a dividere ogni segmento in due parti uguali.
}

\renewcommand{\Risposta}{\ldots}

\renewcommand{\Aiutino}{
Dato un segmento $AB$ conviene costruire un triangolo $ABC$ ed una sua mediana $BM$, trisecando la quale si trova il \ldots\ di $ABC$.
}

\renewcommand{\SoluzioneUfficiale}{
Sia $AB$ il segmento che vogliamo bisecare.
\begin{itemize}
	\item  Scegliamo un punto $P$ che non stia sulla retta $AB$ (Figura 1).

	\item  Usiamo il trisecatore sul segmento $AP$, tracciando 2 punti $N$ e $C$ sul segmento stesso tali che $AN=NC=CP$ (Figura 2).

	\item Usiamo il trisecatore sul segmento $BN$, che è una mediana del triangolo $ABC$: otteniamo altri 2 punti (Figura 3), di cui quello più vicino ad $N$ è il baricentro $G$ di $ABC$ (per la nota proprietà secondo la quale il baricentro divide ogni mediana in 2 parti, di cui quella contenente il vertice è il doppio dell'altra).

	\item Sia $M$ il punto di intersezione tra $CG$ ed $AB$ (Figura 4): dal momento che $CM$ passa per il baricentro di $ABC$, ne segue che $CM$ è una mediana del triangolo, e quindi $M$ è il punto medio di $AB$ (il fatto che la retta $CG$ contenga la mediana assicura in particolare che interseca la retta $AB$).
\end{itemize}

\hfill
\psset{unit=6mm}
\pspicture(-1,-2.5)(5.5,6)
\psline[linewidth=.5\pslinewidth](-1,-1.5)(6,-1.5)(6,5.5)(-1,5.5)(-1,-1.5)
\rput{0}(2.5,-2){Figura 1}

\pstGeonode[PosAngle={270,270,150}](0,0){A}(5,0){B}(1.7,4){P}
\pstLineAB{A}{B}
% \pstHomO[PosAngle=150,HomCoef= 1 3 div]{A}{P}{N}
% \pstHomO[PosAngle=150,HomCoef= 2 3 div]{A}{P}{C}
% 
% \pstHomO[PointName=none,HomCoef= 1 3 div]{B}{N}{GG}
% \pstHomO[PosAngle=70,HomCoef= 2 3 div]{B}{N}{G}
% \pstInterLL[PosAngle=270]{C}{G}{A}{B}{M}
% 
% \psset{linecolor=blue}
% \pstLineAB{A}{C}
% \pstLineAB{B}{C}
% \pstLineAB{B}{N}
% \pstLineAB{M}{C}

\endpspicture
\hfill
\psset{unit=6mm}
\pspicture(-1,-2.5)(5.5,6)
\psline[linewidth=.5\pslinewidth](-1,-1.5)(6,-1.5)(6,5.5)(-1,5.5)(-1,-1.5)
\rput{0}(2.5,-2){Figura 2}

\pstGeonode[PosAngle={270,270,150}](0,0){A}(5,0){B}(1.7,4){P}
\pstLineAB{A}{B}
\pstHomO[PosAngle=150,HomCoef= 1 3 div]{A}{P}[N]
\pstHomO[PosAngle=150,HomCoef= 2 3 div]{A}{P}[C]

% \pstHomO[PointName=none,HomCoef= 1 3 div]{B}{N}{GG}
% \pstHomO[PosAngle=70,HomCoef= 2 3 div]{B}{N}{G}
% \pstInterLL[PosAngle=270]{C}{G}{A}{B}{M}
% 
% \psset{linecolor=blue}
% \pstLineAB{A}{C}
% \pstLineAB{B}{C}
% \pstLineAB{B}{N}
% \pstLineAB{M}{C}

\endpspicture
\hfill
\mbox{}
\bigskip

\hfill
\psset{unit=6mm}
\pspicture(-1,-2.5)(5.5,6)
\psline[linewidth=.5\pslinewidth](-1,-1.5)(6,-1.5)(6,5.5)(-1,5.5)(-1,-1.5)
\rput{0}(2.5,-2){Figura 3}

\pstGeonode[PosAngle={270,270,150}](0,0){A}(5,0){B}(1.7,4){P}
\pstLineAB{A}{B}
\pstHomO[PosAngle=150,HomCoef= 1 3 div]{A}{P}[N]
\pstHomO[PosAngle=150,HomCoef= 2 3 div]{A}{P}[C]

\pstHomO[PointName=none,HomCoef= 1 3 div]{B}{N}[GG]
\pstHomO[PosAngle=100,HomCoef= 2 3 div]{B}{N}[G]
% \pstInterLL[PosAngle=270]{C}{G}{A}{B}{M}

\psset{linecolor=blue}
\pstLineAB{A}{C}
\pstLineAB{B}{C}
\pstLineAB{B}{N}
% \pstLineAB{M}{C}

\endpspicture
\hfill
\psset{unit=6mm}
\pspicture(-1,-2.5)(5.5,6)
\psline[linewidth=.5\pslinewidth](-1,-1.5)(6,-1.5)(6,5.5)(-1,5.5)(-1,-1.5)
\rput{0}(2.5,-2){Figura 4}

\pstGeonode[PosAngle={270,270,150}](0,0){A}(5,0){B}(1.7,4){P}
\pstLineAB{A}{B}
\pstHomO[PosAngle=150,HomCoef= 1 3 div]{A}{P}[N]
\pstHomO[PosAngle=150,HomCoef= 2 3 div]{A}{P}[C]

\pstHomO[PointName=none,HomCoef= 1 3 div]{B}{N}[GG]
\pstHomO[PosAngle=20,HomCoef= 2 3 div]{B}{N}[G]
\pstInterLL[PosAngle=270]{C}{G}{A}{B}{M}

\psset{linecolor=blue}
\pstLineAB{A}{C}
\pstLineAB{B}{C}
\pstLineAB{B}{N}
\pstLineAB{M}{C}

\endpspicture
\hfill
\mbox{}
}

\renewcommand{\SoluzioniAlternative}{
}

\renewcommand{\Osservazioni}{
}

\renewcommand{\Commenti}{
}